% Part 1 of Day 10: Small language models from a systems view.
%
% Topics of the syllabus: parameters, precision, and disk size; quantization
% levels and the GGUF format; RAM budget: weights, KV cache, and the rest of
% the system; tokens per second and what is usable; why distillation makes
% small models good. Exercise: calculate the RAM of a 3-billion-parameter
% model at 4 bits with two context lengths.
%
% The numbers of this part have three origins. Each frame says so in its
% source line.
% 1. The GGUF headers of the two lab models, llama3.2:1b and llama3.2:3b of
%    the Ollama library, read for this course (tools/experiments/day10_gguf.py).
%    The shapes, the types, and the counts come from the files.
% 2. Calculations of this course from these headers and from the block sizes
%    of ggml-common.h of llama.cpp (tools/experiments/day10_slm.py).
% 3. Numbers that the authors of the sources published: the speed on the
%    Raspberry Pi 5 (lab "Small Language Models" of Machine Learning Systems)
%    and the rules of the guide "SLMs at the Edge". Nobody measured a speed
%    on a board for this course.
% Sizes in GB are 10^9 bytes, as in the Ollama library.

\theorypart{Small language models from a systems view}%
  {calculate the RAM of a 4-bit model with 3 billion parameters}

% ------------------------------------------------------------------ the loop
\begin{frame}{A language model writes one token at a time}
  \centering
  \begin{tikzpicture}[font=\small, >=stealth, x=0.9cm,
      box/.style={draw=coolgray, rounded corners=2pt, fill=boxgray,
                  minimum height=0.9cm, align=center}]
    \node[box, text width=2.8cm] (in) at (0,0)
      {Tokens so far\\ \scriptsize prompt + answer};
    \node[box, text width=2.2cm, fill=darkcyan!15] (m) at (4.1,0)
      {Model\\ \scriptsize all weights};
    \node[box, text width=2.6cm] (p) at (8.0,0)
      {Probability of each of 128\,256 tokens};
    \node[box, text width=1.7cm, fill=cardinalred!12] (t) at (8.0,-2.0)
      {Next token};
    \draw[->, thick] (in) -- (m);
    \draw[->, thick] (m) -- (p);
    \draw[->, thick] (p) -- node[right, font=\scriptsize] {select} (t);
    \draw[->, thick] (t) -| node[pos=0.25, below, font=\scriptsize]
      {append, and run the model again} (in);
  \end{tikzpicture}
  \vskip 0.6em
  \small
  \begin{itemize}
    \item A classifier of Days 1 to 9 runs once for one input. A language
          model runs once for each token of its answer.
    \item An answer of 100 tokens needs 100 runs of the model. The speed is
          a rate: tokens in each second.
    \item The model stops at a stop token, or at a limit of the runtime.
  \end{itemize}
  \keypoint{The cost of an answer grows with its length. A short answer is
    a fast answer.}
\end{frame}

% -------------------------------------------------------------------- tokens
\begin{frame}{Tokens}
  \small
  \begin{itemize}
    \item A tokenizer cuts the text into tokens: words, parts of words, and
          signs. The model sees only the numbers of the tokens.
    \item The vocabulary of Llama 3.2 has 128\,256 tokens. The tokenizer
          joins small parts with 280\,147 merge rules (byte-pair encoding).
    \item The GGUF file of the model contains the tokenizer. You need no
          other file.
    \item The context is the number of tokens that the model can see: the
          prompt and the answer. Llama 3.2 is trained for 131\,072 tokens.
  \end{itemize}
  \vskip 0.4em
  \renewcommand{\arraystretch}{1.2}
  \begin{tabular}{@{}L{0.40\textwidth}L{0.50\textwidth}@{}}
    \toprule
    \textbf{Rate of a person} & \textbf{Rate of a model} \\
    \midrule
    Reading: 5 to 7 words in each second & About 8 to 10 tokens in each
      second \\
    \bottomrule
  \end{tabular}
  \vskip 0.4em
  \keypoint{Count in tokens, not in words. The prompt, the context, the
    speed, and the price of a cloud model all use tokens.}
  \source{Vocabulary, merge rules, and context: the GGUF header of
    \code{llama3.2:3b}.\\ Reading rate: guide ``SLMs at the Edge'',
    section 2.6.}
\end{frame}

% ------------------------------------------------------------- two phases
\begin{frame}{Two phases: the prompt and the generation}
  \small
  \begin{columns}[T]
    \begin{column}{0.48\textwidth}
      \textbf{Prompt evaluation (prefill)}
      \begin{itemize}
        \item All tokens of the prompt go through the model together.
        \item The weights multiply a matrix of tokens: many operations for
              each byte that the processor reads.
        \item The result: the time to the first token.
      \end{itemize}
    \end{column}
    \begin{column}{0.48\textwidth}
      \textbf{Generation (decode)}
      \begin{itemize}
        \item One new token for each run of the model.
        \item The weights multiply one vector: few operations for each
              byte.
        \item The result: the tokens in each second of the answer.
      \end{itemize}
    \end{column}
  \end{columns}
  \vskip 0.5em
  \footnotesize
  \renewcommand{\arraystretch}{1.15}
  \begin{tabular}{@{}L{0.42\textwidth}rrr@{}}
    \toprule
    \textbf{\code{llama3.2:1b}, Raspberry Pi 5} & \textbf{Tokens}
    & \textbf{Time} & \textbf{Rate} \\
    \midrule
    Prompt eval  & 32  & 1.64 s  & 19.46 tokens/s \\
    Eval (the answer) & 8 & 0.89 s & 8.99 tokens/s \\
    \bottomrule
  \end{tabular}
  \vskip 0.3em
  \small
  Ollama prints the two rates with \code{--verbose}. A benchmark of a
  language model gives the two rates and the two token counts.
  \source{Lab ``Small Language Models'' of \emph{Machine Learning Systems}.
    \\ The token count of the answer is $8.99 \times 0.89$ s, a
    calculation.}
\end{frame}

% ----------------------------------------------------------- inside a model
\begin{frame}{Inside the two lab models}
  \small
  \renewcommand{\arraystretch}{1.15}
  \begin{tabular}{@{}L{0.48\textwidth}rr@{}}
    \toprule
    \textbf{Value in the GGUF header} & \code{llama3.2:1b} & \code{llama3.2:3b} \\
    \midrule
    Layers (transformer blocks) & 16      & 28 \\
    Width of the model, $d$     & 2048    & 3072 \\
    Attention heads             & 32      & 24 \\
    Key and value heads         & 8       & 8 \\
    Size of one head            & 64      & 128 \\
    Width of the MLP            & 8192    & 8192 \\
    Vocabulary                  & 128\,256 & 128\,256 \\
    Parameters                  & 1\,235\,814\,432 & 3\,212\,749\,888 \\
    \bottomrule
  \end{tabular}
  \vskip 0.4em
  \begin{itemize}
    \item The names say 1B and 3B. The files have 1.24 and 3.21 billion
          parameters.
    \item Fewer key and value heads than attention heads: grouped-query
          attention. It makes the KV cache smaller (later in this part).
  \end{itemize}
  \source{The GGUF headers of the two files of the Ollama library, read for
    this course.}
\end{frame}

% --------------------------------------------------------- count parameters
\begin{frame}{Count the parameters of the 3B model}
  \small
  \renewcommand{\arraystretch}{1.15}
  \begin{tabular}{@{}L{0.30\textwidth}L{0.40\textwidth}r@{}}
    \toprule
    \textbf{Part} & \textbf{Calculation} & \textbf{Parameters} \\
    \midrule
    Attention: Q and output & $2 \times 3072 \times 3072$ & 18\,874\,368 \\
    Attention: K and V      & $2 \times 3072 \times (8 \times 128)$ & 6\,291\,456 \\
    MLP: gate, up, down     & $3 \times 3072 \times 8192$ & 75\,497\,472 \\
    Two norms               & $2 \times 3072$ & 6\,144 \\
    \midrule
    One layer               & sum of the four lines & 100\,669\,440 \\
    28 layers               & $28 \times 100\,669\,440$ & 2\,818\,744\,320 \\
    Token table             & $128\,256 \times 3072$ & 394\,002\,432 \\
    Last norm and RoPE      & $3072 + 64$ & 3\,136 \\
    \midrule
    \textbf{Total}          & & \textbf{3\,212\,749\,888} \\
    \bottomrule
  \end{tabular}
  \vskip 0.4em
  \begin{itemize}
    \item The MLP has 75 percent of the weights of a layer.
    \item The token table is also the output layer (tied weights). It has
          12 percent of all parameters.
  \end{itemize}
  \source{Shapes: the GGUF header of \code{llama3.2:3b}.\\ The total is the
    sum of the tensors in the file.}
\end{frame}

% ------------------------------------------------------------- disk size
\begin{frame}{Precision and size}
  \small
  \textbf{Size in bytes = parameters $\times$ bits for each weight / 8}
  \vskip 0.4em
  \renewcommand{\arraystretch}{1.12}
  \begin{tabular}{@{}L{0.30\textwidth}rrr@{}}
    \toprule
    \textbf{Format} & \textbf{Bits} & \textbf{1.24 B parameters} & \textbf{3.21 B parameters} \\
    \midrule
    \code{float32}            & 32   & 4.94 GB & 12.85 GB \\
    \code{float16}            & 16   & 2.47 GB & 6.43 GB \\
    \code{Q8\_0}              & 8.5  & 1.31 GB & 3.41 GB \\
    \code{Q4\_0}              & 4.5  & 0.70 GB & 1.81 GB \\
    Only the 4 bits           & 4    & 0.62 GB & 1.61 GB \\
    \midrule
    The file of the lab       &      & 1.32 GB, \code{Q8\_0} & 2.02 GB, \code{Q4\_K\_M} \\
    \bottomrule
  \end{tabular}
  \vskip 0.4em
  \begin{itemize}
    \item Training gives \code{float16} or \code{float32}. A 3B model in
          \code{float16} uses most of the 8 GB of the Raspberry Pi 5.
    \item A 4-bit format needs more than 4 bits: each block of weights also
          stores its scale.
  \end{itemize}
  \keypoint{Quantization makes a small language model fit on the board.}
  \source{Parameters: the GGUF headers.\\ File sizes: the Ollama registry.}
\end{frame}

% ------------------------------------------------------------- blocks
\begin{frame}{Quantization in blocks: \code{Q8\_0} and \code{Q4\_0}}
  \centering
  \begin{tikzpicture}[font=\scriptsize, x=0.33cm, y=0.5cm]
    \fill[cardinalred!25] (0,0) rectangle (2,1);
    \node at (1,0.5) {$d$};
    \foreach \i in {0,...,31} {
      \draw[coolgray, fill=darkcyan!12] (2+\i*0.75,0) rectangle (2.75+\i*0.75,1);
    }
    \draw[coolgray] (0,0) rectangle (2,1);
    \node[anchor=west] at (0,1.6) {\code{Q8\_0}: one block of 32 weights = 2 + 32 bytes = 34 bytes};
    \draw[coolgray, |-|] (2,-0.15) -- (26,-0.15)
      node[midway, below=1pt] {32 integers $q$ of 8 bits};
    \node[anchor=north] at (1,-0.2) {FP16 scale};
    \begin{scope}[yshift=-2.7cm]
      \fill[cardinalred!25] (0,0) rectangle (2,1);
      \draw[coolgray] (0,0) rectangle (2,1);
      \node at (1,0.5) {$d$};
      \foreach \i in {0,...,15} {
        \draw[coolgray, fill=darkcyan!12] (2+\i*0.75,0) rectangle (2.75+\i*0.75,1);
      }
      \node[anchor=west] at (0,1.6) {\code{Q4\_0}: one block of 32 weights = 2 + 16 bytes = 18 bytes};
      \draw[coolgray, |-|] (2,-0.15) -- (14,-0.15)
        node[midway, below=1pt] {32 integers of 4 bits, two in each byte};
    \end{scope}
  \end{tikzpicture}
  \vskip 0.5em
  \small
  \begin{itemize}
    \item Each weight is $w = d \times q$. Each block of 32 weights has its
          own scale $d$.
    \item \code{Q8\_0}: $34 \times 8 / 32 = 8.5$ bits for each weight.
          \code{Q4\_0}: $18 \times 8 / 32 = 4.5$ bits.
    \item Day 6 used one scale for each tensor or for each channel. A small
          block gives a better fit to the weights, for some more bytes.
  \end{itemize}
  \source{Block definitions in \code{ggml-common.h} of llama.cpp (MIT).}
\end{frame}

% ------------------------------------------------------------- k-quants
\begin{frame}{K-quants: blocks in a super-block}
  \small
  A K-quant groups 256 weights into a super-block of 8 blocks with 32
  weights. Each block has a scale and a minimum of 6 bits. The super-block
  has two FP16 values for these scales: $w = a \times q + b$.
  \vskip 0.4em
  \renewcommand{\arraystretch}{1.15}
  \begin{tabular}{@{}lL{0.42\textwidth}rr@{}}
    \toprule
    \textbf{Type} & \textbf{Bytes of one super-block} & \textbf{Bytes} &
    \textbf{Bits for each weight} \\
    \midrule
    \code{Q4\_K} & $2 + 2 + 12 + 128$             & 144 & 4.5 \\
    \code{Q5\_K} & $2 + 2 + 12 + 32 + 128$        & 176 & 5.5 \\
    \code{Q6\_K} & $2 + 16 + 64 + 128$            & 210 & 6.5625 \\
    \bottomrule
  \end{tabular}
  \vskip 0.4em
  \begin{itemize}
    \item The bits per weight: bytes $\times$ 8 / 256. Example:
          $144 \times 8 / 256 = 4.5$.
    \item The minimum $b$ gives a better fit when the weights of a block
          are not centred on zero.
    \item The scales themselves are quantized: 6 bits, not 16. So the cost
          of 8 scales stays small.
  \end{itemize}
  \source{Block definitions in \code{ggml-common.h} of llama.cpp (MIT).}
\end{frame}

% ------------------------------------------------------------- names
\begin{frame}{Read the name of a quantization level}
  \centering
  \begin{tikzpicture}[font=\small, >=stealth]
    \node[font=\LARGE\ttfamily] (n) at (0,0) {Q4\_K\_M};
    \node[align=center, text width=3.0cm, font=\scriptsize] (a) at (-3.6,-1.5)
      {\textbf{4}: the main number of bits};
    \node[align=center, text width=3.0cm, font=\scriptsize] (b) at (0,-1.7)
      {\textbf{K}: a K-quant with super-blocks};
    \node[align=center, text width=3.4cm, font=\scriptsize] (c) at (3.8,-1.5)
      {\textbf{M}: medium mix. Some tensors keep more bits. S: small, L:
      large.};
    \draw[->, coolgray] (-0.55,-0.3) -- (a.north);
    \draw[->, coolgray] (0.05,-0.3) -- (b.north);
    \draw[->, coolgray] (0.65,-0.3) -- (c.north);
  \end{tikzpicture}
  \vskip 0.6em
  \small
  \renewcommand{\arraystretch}{1.15}
  \begin{tabular}{@{}L{0.18\textwidth}L{0.74\textwidth}@{}}
    \toprule
    \textbf{Level} & \textbf{Advice of the guide} \\
    \midrule
    \code{Q8\_0}   & Almost no loss. For models of 1B or smaller. \\
    \code{Q5\_K\_M} & A good compromise near 1B, for example for
                     structured output. \\
    \code{Q4\_K\_M} & The default of Ollama. For models of 2B and larger. \\
    \code{Q3}, \code{Q2} & The quality drops fast. A 1B model at
                     \code{Q2} is often worse than a 360M model at \code{Q8}. \\
    \bottomrule
  \end{tabular}
  \vskip 0.3em
  \begin{itemize}
    \item A small model loses more quality from fewer bits than a large
          model.
  \end{itemize}
  \source{Guide ``SLMs at the Edge'', section 2.2.\\ Published advice, not
    measurements of this course.}
\end{frame}

% ------------------------------------------------------------- real mix
\begin{frame}{\code{Q4\_K\_M} is a mix: the file of the lab}
  \small
  The tensors of \code{llama3.2:3b}, from its GGUF header:
  \vskip 0.3em
  \footnotesize
  \renewcommand{\arraystretch}{1.15}
  \begin{tabular}{@{}lrL{0.38\textwidth}rr@{}}
    \toprule
    \textbf{Type} & \textbf{Tensors} & \textbf{Which tensors} &
    \textbf{Weights} & \textbf{Bytes} \\
    \midrule
    \code{Q4\_K} & 168 & Q, K, output, gate, up; one half of V and down
                 & 2.422 B & 1.362 GB \\
    \code{Q6\_K} & 29  & The token table; the other half of V and down
                 & 0.790 B & 0.648 GB \\
    \code{F32}   & 58  & The norms and the RoPE values
                 & 0.0002 B & 0.0007 GB \\
    \midrule
    Sum          & 255 & & 3.213 B & 2.012 GB \\
    \bottomrule
  \end{tabular}
  \vskip 0.4em
  \small
  \begin{itemize}
    \item Mean: $2.012 \times 8 / 3.213 = 5.0$ bits for each weight, not 4.
    \item The tensors V and down of 14 of the 28 layers keep 6 bits. These
          tensors are the most sensitive to the error.
    \item The file has 2.019 GB. The other 7.8 MB are the header: most of
          it is the tokenizer.
  \end{itemize}
  \keypoint{Calculate the RAM from the size of the file, not from the name of
    the level.}
  \source{The GGUF header of \code{llama3.2:3b}.}
\end{frame}

% ------------------------------------------------------------- GGUF
\begin{frame}[fragile]{The GGUF file format}
  \begin{columns}[T]
    \begin{column}{0.40\textwidth}
      \centering
      \begin{tikzpicture}[font=\scriptsize,
          seg/.style={draw=coolgray, minimum width=3.6cm, align=center}]
        \node[seg, fill=cardinalred!15, minimum height=0.6cm] (h) at (0,0)
          {Header: \code{GGUF}, version 3,\\ counts};
        \node[seg, fill=boxgray, minimum height=1.0cm, below=0pt of h] (m)
          {Metadata: 30 keys\\ architecture, shapes,\\ tokenizer, chat template};
        \node[seg, fill=boxgray, minimum height=0.7cm, below=0pt of m] (t)
          {Tensor list: 255 names,\\ shapes, types, offsets};
        \node[seg, fill=darkcyan!15, minimum height=2.2cm, below=0pt of t] (d)
          {Tensor data\\ 2.012 GB};
      \end{tikzpicture}
    \end{column}
    \begin{column}{0.57\textwidth}
\begin{lstlisting}
general.architecture   llama
general.name           Llama 3.2 3B Instruct
general.file_type      15 (Q4_K_M)
llama.block_count      28
llama.embedding_length 3072
llama.context_length   131072
llama.attention.head_count_kv  8
tokenizer.ggml.tokens  128256 strings
\end{lstlisting}
      \small
      \begin{itemize}
        \item One file has the weights, the tokenizer, and the chat
              template.
        \item llama.cpp, Ollama, and LM Studio read GGUF.
        \item The runtime can map the tensor data into memory: it loads
              only the pages that it reads.
      \end{itemize}
    \end{column}
  \end{columns}
  \source{Keys and counts: the GGUF header of \code{llama3.2:3b}, read for
    this course.}
\end{frame}

% ------------------------------------------------------------- RAM budget
\begin{frame}{The RAM budget}
  \centering
  \begin{tikzpicture}[font=\scriptsize, x=1.25cm, y=0.8cm]
    \fill[coolgray!35]     (0,0)   rectangle (0.4,1);
    \fill[darkcyan!35]     (0.4,0) rectangle (2.42,1);
    \fill[cardinalred!30]  (2.42,0) rectangle (2.89,1);
    \fill[darkbeige!40]    (2.89,0) rectangle (3.4,1);
    \draw[coolgray] (0,0) rectangle (8,1);
    \node[rotate=90] at (0.2,0.5) {OS};
    \node at (1.41,0.5) {weights 2.02 GB};
    \node[rotate=90, font=\tiny] at (2.655,0.5) {KV};
    \node[rotate=90, font=\tiny] at (3.145,0.5) {other};
    \node at (5.7,0.5) {free for your application};
    \foreach \x in {0,2,4,6,8} { \node[below] at (\x,0) {\x\ GB}; }
  \end{tikzpicture}
  \vskip 0.3em
  {\scriptsize Example: \code{llama3.2:3b} with a context of 4096 tokens on a
  Raspberry Pi 5 with 8 GB. The width of ``other'' is not a measurement.\par}
  \vskip 0.4em
  \small
  \begin{enumerate}
    \item \textbf{Weights:} the size of the file.
    \item \textbf{KV cache:} it grows with the context length.
    \item \textbf{Other:} buffers of the runtime, the operating system, and
          your program.
  \end{enumerate}
  \begin{itemize}
    \item The source measured 3.24 GB in use with a 3.8B model loaded,
          and about 377 MB after Ollama stopped (no desktop).
    \item Rule of the guide: usable model size $\approx$ (total RAM $-$
          operating system and applications) $\times$ 0.7. For 8 GB:
          $(8 - 0.4) \times 0.7 = 5.3$ GB for the weights and the cache.
  \end{itemize}
  \source{Memory values: lab ``Small Language Models'' of \emph{Machine
    Learning Systems}.\\ Rule: guide ``SLMs at the Edge'', section 2.4.}
\end{frame}

% ------------------------------------------------------------- KV cache
\begin{frame}{The KV cache}
  \begin{columns}[c]
    \begin{column}{0.45\textwidth}
      \centering
      \begin{tikzpicture}[font=\scriptsize, x=0.5cm, y=0.5cm]
        \foreach \i in {0,...,6} {
          \draw[coolgray, fill=darkcyan!15] (\i,2) rectangle (\i+1,3);
          \draw[coolgray, fill=cardinalred!12] (\i,0.6) rectangle (\i+1,1.6);
        }
        \draw[coolgray, fill=darkcyan!45] (7,2) rectangle (8,3);
        \draw[coolgray, fill=cardinalred!40] (7,0.6) rectangle (8,1.6);
        \node[anchor=east] at (-0.1,2.5) {K};
        \node[anchor=east] at (-0.1,1.1) {V};
        \node[anchor=north] at (3.5,0.4) {tokens before: stored};
        \node[anchor=south] at (7.5,3.1) {new};
      \end{tikzpicture}
    \end{column}
    \begin{column}{0.53\textwidth}
      \small
      \begin{itemize}
        \item Attention compares the new token with all tokens before it.
              It needs their keys K and values V in each layer.
        \item The runtime stores K and V. It does not calculate them again
              for each new token.
        \item The cache grows by one row for each token.
      \end{itemize}
    \end{column}
  \end{columns}
  \vskip 0.6em
  \centering
  \textbf{Bytes = 2 $\times$ layers $\times$ KV heads $\times$ head size
  $\times$ context $\times$ bytes for each value}
  \vskip 0.4em
  \raggedright
  \small
  \begin{itemize}
    \item The 2 counts K and V. Ollama stores the values in \code{f16}: 2
          bytes. \code{q8\_0} uses about one half of this memory.
    \item The default context of Ollama is 4096 tokens. The runtime
          allocates the cache for this context.
  \end{itemize}
  \source{Ollama FAQ: the context window and the type of the K/V cache.}
\end{frame}

\begin{frame}{The KV cache of the two lab models}
  \small
  \renewcommand{\arraystretch}{1.15}
  \begin{tabular}{@{}L{0.36\textwidth}rr@{}}
    \toprule
    & \code{llama3.2:1b} & \code{llama3.2:3b} \\
    \midrule
    Calculation for one token
      & $2 \cdot 16 \cdot 8 \cdot 64 \cdot 2$ & $2 \cdot 28 \cdot 8 \cdot 128 \cdot 2$ \\
    Bytes for one token         & 32\,768  & 114\,688 \\
    4096 tokens (default)       & 0.13 GB  & 0.47 GB \\
    32\,768 tokens              & 1.07 GB  & 3.76 GB \\
    131\,072 tokens (maximum)   & 4.29 GB  & 15.03 GB \\
    \midrule
    The weights (file)          & 1.32 GB  & 2.02 GB \\
    \bottomrule
  \end{tabular}
  \vskip 0.4em
  \begin{itemize}
    \item With the maximum context, the cache of the 3B model is 7 times
          its weights. It does not fit in 8 GB.
    \item With 24 KV heads instead of 8, the 3B cache would be 3 times as
          large: grouped-query attention saves this memory.
  \end{itemize}
  \keypoint{Set the context to the length that your application needs, not
    to the maximum of the model.}
  \source{Calculation of this course from the GGUF headers,\\ with
    \code{f16} values.}
\end{frame}

% ------------------------------------------------------------- memory-bound
\begin{frame}{Generation is memory-bound}
  \small
  For each new token, the model reads each weight once and does about 2
  operations with it.
  \vskip 0.4em
  \renewcommand{\arraystretch}{1.15}
  \begin{tabular}{@{}L{0.46\textwidth}L{0.44\textwidth}@{}}
    \toprule
    \textbf{\code{llama3.2:3b}, one token} & \textbf{Value} \\
    \midrule
    Operations                  & $2 \times 3.21$ B = 6.4 GFLOP \\
    Bytes read                  & 2.02 GB (the file) \\
    Intensity                   & $6.4 / 2.02 = 3.2$ FLOP for each byte \\
    At 5.3 tokens/s (published) & $6.4 \times 5.3 = 34$ GFLOP/s \\
    \bottomrule
  \end{tabular}
  \vskip 0.4em
  \begin{itemize}
    \item Day 7: the peak of the 4 cores is 153.6 GFLOP/s. The generation
          uses less than one quarter of it.
    \item An intensity of 3.2 is low: the memory sets the limit, as for the
          matrix times vector of Day 7.
    \item The prompt is different: all its tokens use each weight. Its
          intensity is higher, so its rate is higher.
  \end{itemize}
  \source{Calculation of this course. Rate: lab ``Small Language Models''
    of \emph{Machine Learning Systems}.}
\end{frame}

\begin{frame}{The published rates agree with the memory limit}
  \small
  \renewcommand{\arraystretch}{1.15}
  \begin{tabular}{@{}L{0.22\textwidth}rrrr@{}}
    \toprule
    \textbf{Model} & \textbf{Parameters} & \textbf{File} & \textbf{Eval rate}
    & \textbf{File $\times$ rate} \\
    \midrule
    \code{llama3.2:1b} & 1.24 B & 1.32 GB & 8.99 tokens/s & 11.9 GB/s \\
    \code{llama3.2:3b} & 3.21 B & 2.02 GB & 5.3 tokens/s  & 10.7 GB/s \\
    \midrule
    Ratio 3B / 1B      & 2.6    & 1.53    & 1.70 (slower) & \\
    \bottomrule
  \end{tabular}
  \vskip 0.4em
  \begin{itemize}
    \item The 3B model has 2.6 times the parameters, but it is only 1.7
          times slower. Its file is only 1.53 times as large.
    \item The product of the file size and the rate is almost the same for
          the two models: about 11 GB of weights in each second.
    \item So the bytes set the speed, not the parameters. Fewer bits give
          more tokens in each second.
  \end{itemize}
  \keypoint{Upper limit: tokens per second $\approx$ memory bandwidth / size
    of the file. Measure the rate in the lab to see your limit.}
  \source{Rates: lab ``Small Language Models'' of \emph{Machine Learning
    Systems}, about 2 GB for the\\ 3B file in the source. Products and
    ratios: calculations of this course.}
\end{frame}

% ------------------------------------------------------------- usable
\begin{frame}{Tokens per second: what is usable?}
  \small
  \renewcommand{\arraystretch}{1.2}
  \begin{tabular}{@{}L{0.30\textwidth}L{0.24\textwidth}L{0.36\textwidth}@{}}
    \toprule
    \textbf{Use} & \textbf{Rate that you need} & \textbf{Why} \\
    \midrule
    Chat with a person & 8 to 10 tokens/s & A person reads at this rate \\
    A job with no person: summarize a log, classify a message
      & 2 tokens/s can be enough & Nobody waits for each word \\
    \bottomrule
  \end{tabular}
  \vskip 0.5em
  \textbf{The total time of one answer} (if the rates stay the same):
  \vskip 0.2em
  \begin{itemize}
    \item A prompt of 500 tokens at 19.46 tokens/s: 25.7 s before the
          first token.
    \item An answer of 100 tokens: 11.1 s with the 1B model, 18.9 s with
          the 3B model.
    \item A reasoning model first writes a ``think'' text. The person waits
          for these tokens too.
  \end{itemize}
  \keypoint{The time to the first token depends on the prompt. A long
    prompt can cost more time than the answer.}
  \source{Uses and rates of a person: guide ``SLMs at the Edge'', section
    2.6. Times: calculations\\ of this course with the published rates.}
\end{frame}

% ------------------------------------------------------------- distillation
\begin{frame}{Why small models got good: distillation}
  \begin{columns}[c]
    \begin{column}{0.45\textwidth}
      \centering
      \begin{tikzpicture}[font=\scriptsize, >=stealth,
          box/.style={draw=coolgray, rounded corners=2pt, align=center}]
        \node[box, fill=darkcyan!15, minimum width=2.6cm, minimum height=1.3cm]
          (T) at (0,0) {Teacher\\ large model};
        \node[box, fill=cardinalred!12, minimum width=1.6cm, minimum height=0.8cm]
          (S) at (0,-2.6) {Student\\ small model};
        \node[box, fill=boxgray] (X) at (-2.6,-1.3) {Text};
        \draw[->] (X) |- (T);
        \draw[->] (X) |- (S);
        \draw[->, thick] (T) -- node[right, align=left]
          {probability of\\ each next token} (S);
      \end{tikzpicture}
    \end{column}
    \begin{column}{0.53\textwidth}
      \small
      \begin{itemize}
        \item The student learns the full probabilities of the teacher, not
              only the correct next token.
        \item These soft targets also show which wrong tokens are almost
              correct. This is more information for each example.
        \item Day 6 used the same idea for a classifier.
      \end{itemize}
    \end{column}
  \end{columns}
  \vskip 0.5em
  \small
  \textbf{Llama 3.2 1B and 3B:} the authors pruned the 8B model of Llama
  3.1. Then the outputs of the 8B and 70B models were the targets for each
  token. The training used 9 trillion tokens.
  \keypoint{A small model with good data and a good teacher can do a narrow
    task well. It does not have the breadth of the teacher.}
  \source{Hinton, Vinyals, and Dean, 2015. Llama 3.2: lab ``Small Language
    Models'' of\\ \emph{Machine Learning Systems}.}
\end{frame}

% ------------------------------------------------------------- scope
\begin{frame}{What a small model does well}
  \small
  \renewcommand{\arraystretch}{1.2}
  \begin{tabular}{@{}L{0.45\textwidth}L{0.45\textwidth}@{}}
    \toprule
    \textbf{Good tasks for a model of about 1B} & \textbf{Weak tasks} \\
    \midrule
    Summarize a text                     & Reasoning in many steps \\
    Find the intent of a short command   & Facts that are not in the prompt \\
    Answer questions on a narrow topic   & Exact arithmetic \\
    Extract fields from messy text       & Long and complex instructions \\
    A basic conversation                 & A guaranteed correct answer \\
    \bottomrule
  \end{tabular}
  \vskip 0.5em
  \begin{itemize}
    \item The guide: the question at the edge is not ``how large a model
          fits'', but ``how small a model still does the task''.
    \item Part 3 adds tools that help a small model: a schema for the
          output, function calls, and retrieval.
  \end{itemize}
  \keypoint{Test the model with your own prompts. A ranking on a public
    leaderboard does not tell you the result on your task.}
  \source{Guide ``SLMs at the Edge'', sections 1 and 7.}
\end{frame}

% -------------------------------------------------------------------- exercise
\exerciseframe{the RAM of a 4-bit 3B model}{%
  \small
  A model has 3.2 billion parameters, 28 layers, 8 key and value heads, and
  a head size of 128. The runtime stores the KV cache in \code{f16}. The
  board is a Raspberry Pi 5 with 8 GB. The operating system uses 0.4 GB.
  \begin{enumerate}
    \item Calculate the size of the weights at 4 bits for each weight.
    \item Calculate the bytes of the KV cache for one token.
    \item Calculate the RAM of the weights and the KV cache for a context of
          4096 tokens and for a context of 32\,768 tokens.
    \item Use the rule of the guide: (8 $-$ 0.4) $\times$ 0.7 = 5.3 GB. Which
          context fits?
    \item The file of the lab has 2.02 GB, not your value of question 1.
          Why?
    \item Give two changes that make the larger context fit.
  \end{enumerate}
  Work with your neighbour for 8 minutes. Use GB = $10^9$ bytes.
}

\answerframe{the RAM of a 4-bit 3B model}{%
  \small
  \begin{enumerate}
    \item $3.2 \times 10^9 \times 4 / 8 = 1.6$ GB.
    \item $2 \times 28 \times 8 \times 128 \times 2 = 114\,688$ bytes.
    \item 4096 tokens: $114\,688 \times 4096 = 0.47$ GB, total
          $1.6 + 0.47 = 2.07$ GB.\\
          32\,768 tokens: $114\,688 \times 32\,768 = 3.76$ GB, total
          $1.6 + 3.76 = 5.36$ GB.
    \item 4096 tokens fits with a large margin. 32\,768 tokens is above
          5.3 GB: it does not fit with the rule. The real file is larger
          too.
    \item Each block stores its scales, and \code{Q4\_K\_M} keeps some
          tensors at 6 bits: 5.0 bits for each weight on average.
    \item A \code{q8\_0} KV cache: 1.88 GB, total 3.48 GB. Or a shorter
          context, for example 16\,384 tokens. Or a smaller model.
  \end{enumerate}
  \keypoint{The weights are fixed. The context is your choice, and it can
    cost more RAM than the weights.}
}
